\section{Falsifiable Predictions and Discriminative Verification Templates}\label{sec:falsifiable}

\subsection{Fibonacci Recurrence Scale Fingerprint}

Under golden-branch scanning, statistical quantities (e.g., autocorrelation of interference visibility, autocorrelation of local patch counting, return time distribution of stabilized sequence) should exhibit Fibonacci-scale stratification or spectral peaks. Verification steps:
\begin{enumerate}
  \item Fix the readout kernel and apparatus parameters, sampling long sequences;
  \item Estimate return-time histograms and autocorrelations;
  \item Compare with periodic approximant (rational rotation numbers) and random control sequences;
  \item Fit peak structures and scale-invariance errors on Fibonacci indices.
\end{enumerate}

\subsection{Kernel Width \texorpdfstring{$\varepsilon$}{epsilon} Control of Interference Visibility Curve}

From the coherent superposition structure in Section \ref{sec:cut-project-readout}, interference visibility should vary continuously with $\varepsilon$: when $\varepsilon$ is small enough to separate internal-coordinate channels, cross-terms are suppressed; when $\varepsilon$ is large enough, cross-terms enhance. The verification strategy is to systematically scan $\varepsilon$ in the same geometric configuration and fit a single-parameter family curve, thereby discriminating between ``kernel-induced complementarity'' and ``ontological bistability-switching'' narratives.

\subsection{Response Kernel Non-Locality and Correlation Bounds}

By changing the apparatus response function (equivalently changing the forward-support weight of $G_{\theta}(\tau)$), the intensity of measured ``apparent backward correlation'' should co-vary with $\Gamma_+$ and significantly decay as the kernel becomes localized. This verification serves to discriminate between ``protocol non-locality'' and ``ontological backward causality'' narratives.

\subsection{Degeneracy Spectrum Frequency Law (Discrete Prototype Exhaustively Verifiable)}

In the discrete prototype, if the microscopic marginal distribution is near-uniform, then the stable-state frequency should approximate $q(x)=d(x)/2^m$. This verification can be directly verified in numerical simulations of reversible permutation dynamics, and serves as a baseline for more complex models.
